\documentclass[11pt,letterpaper]{article}
\usepackage[utf8]{inputenc}
\usepackage[T1]{fontenc}
\usepackage{lmodern}
\usepackage{textcomp}
\usepackage[spanish,es-noshorthands]{babel}
\usepackage[letterpaper,top=3cm,bottom=2.5cm,left=2.5cm,right=2.5cm,headheight=1.1cm,headsep=0.6cm]{geometry}
\usepackage{graphicx}
\usepackage{float}
\usepackage{fancyhdr}
\usepackage{listings}
\usepackage{xcolor}
\usepackage[colorlinks=true,linkcolor=blue,urlcolor=blue]{hyperref}
\urlstyle{same}

\pagestyle{fancy}
\fancyhf{}
\fancyhead[L]{\includegraphics[height=0.8cm]{imagenes/logo-usm.png}}
\renewcommand{\headrulewidth}{0.4pt}
\renewcommand{\footrulewidth}{0.4pt}
\fancyfoot[L]{SIG}
\fancyfoot[C]{\textbf{\thepage}}
\fancyfoot[R]{Lorenzo Ruiz}

\lstset{
  basicstyle=\ttfamily\small,
  columns=fullflexible,
  keepspaces=true,
  breaklines=true,
  upquote=true,
  aboveskip=6pt,
  belowskip=4pt,
  literate={á}{{\'a}}1 {é}{{\'e}}1 {í}{{\'i}}1 {ó}{{\'o}}1 {ú}{{\'u}}1 {ñ}{{\~n}}1
           {Á}{{\'A}}1 {É}{{\'E}}1 {Í}{{\'I}}1 {Ó}{{\'O}}1 {Ú}{{\'U}}1 {Ñ}{{\~N}}1
}

\begin{document}

\begin{titlepage}
\thispagestyle{empty}
\noindent\includegraphics[width=0.39\textwidth]{imagenes/logo-industrias.png}

\vspace*{4.2cm}
\noindent\rule{\textwidth}{0.6pt}
\begin{center}
\vspace{-4pt}
{\Huge\textsc{Laboratorio 5}}\\[2pt]
{\large\textsc{Entregable SQL (2/2): Andes Bike}}
\vspace{-4pt}
\end{center}
\noindent\rule{\textwidth}{0.6pt}

\vspace{2.6cm}
\begin{center}
\textbf{Lorenzo Ruiz}\\
Ingeniería Civil Industrial, 202304619-1, \textit{lflandez@usm.cl}\\
Repositorio: \url{https://github.com/LorenzoRF23/ICN292-Lab5-Ruiz-Lorenzo}
\end{center}

\vfill
\begin{center}
\textsc{noviembre, 2026}\\[0.6cm]
\textsc{Campus Vitacura}
\end{center}
\vspace{1.5cm}
\end{titlepage}

\section*{Consulta 1}
\begin{figure}[H]
\begin{lstlisting}
-- Consulta 1: vendedores con su cantidad de ventas (incluye los que tienen cero)
SELECT v.nombre,
       COUNT(ve.id_venta) AS cantidad_ventas
FROM vendedores v
LEFT JOIN ventas ve ON ve.id_vendedor = v.id_vendedor
GROUP BY v.id_vendedor, v.nombre
ORDER BY cantidad_ventas DESC, v.nombre;
\end{lstlisting}
\centering
\includegraphics[width=\textwidth]{imagenes/c01.png}
\caption{Resultado de la consulta 1}
\end{figure}


\clearpage
\section*{Consulta 2}
\begin{figure}[H]
\begin{lstlisting}
-- Consulta 2: ventas y unidades por tienda (incluye tiendas sin vendedores)
SELECT t.nombre AS tienda,
       COUNT(ve.id_venta)            AS cantidad_ventas,
       COALESCE(SUM(ve.unidades), 0) AS unidades_vendidas
FROM tiendas t
LEFT JOIN vendedores v ON v.id_tienda   = t.id_tienda
LEFT JOIN ventas ve    ON ve.id_vendedor = v.id_vendedor
GROUP BY t.id_tienda, t.nombre
ORDER BY unidades_vendidas DESC, t.nombre;
\end{lstlisting}
\centering
\includegraphics[width=\textwidth]{imagenes/c02.png}
\caption{Resultado de la consulta 2}
\end{figure}


\clearpage
\section*{Consulta 3}
\begin{figure}[H]
\begin{lstlisting}
-- Consulta 3: tiendas sin ningún vendedor asignado
SELECT t.id_tienda, t.nombre, t.comuna
FROM tiendas t
LEFT JOIN vendedores v ON v.id_tienda = t.id_tienda
WHERE v.id_vendedor IS NULL;
\end{lstlisting}
\centering
\includegraphics[width=\textwidth]{imagenes/c03.png}
\caption{Resultado de la consulta 3}
\end{figure}


\clearpage
\section*{Consulta 4}
\begin{figure}[H]
\begin{lstlisting}
-- Consulta 4: modelos que no se vendieron ni una vez en junio de 2025
SELECT b.id_bicicleta, b.modelo, b.categoria
FROM bicicletas b
LEFT JOIN ventas ve
       ON ve.id_bicicleta = b.id_bicicleta
      AND ve.fecha >= '2025-06-01'
      AND ve.fecha <  '2025-07-01'
WHERE ve.id_venta IS NULL
ORDER BY b.id_bicicleta;
\end{lstlisting}
\centering
\includegraphics[width=\textwidth]{imagenes/c04.png}
\caption{Resultado de la consulta 4}
\end{figure}

\noindent\textbf{Explicación:} Puse la condición del mes en el \texttt{ON} del \texttt{LEFT JOIN}, para que solo se unan las ventas de junio y los modelos sin ventas ese mes queden con NULL; si la hubiera puesto en el \texttt{WHERE}, el filtro de fecha habría descartado esas filas con NULL y la consulta habría devuelto un resultado vacío.

\clearpage
\section*{Consulta 5}
\begin{figure}[H]
\begin{lstlisting}
-- Consulta 5: vendedores sin ventas, con el nombre de su tienda
SELECT v.nombre AS vendedor,
       t.nombre AS tienda
FROM vendedores v
LEFT JOIN tiendas t ON t.id_tienda  = v.id_tienda
LEFT JOIN ventas ve ON ve.id_vendedor = v.id_vendedor
WHERE ve.id_venta IS NULL
ORDER BY v.nombre;
\end{lstlisting}
\centering
\includegraphics[width=\textwidth]{imagenes/c05.png}
\caption{Resultado de la consulta 5}
\end{figure}


\clearpage
\section*{Consulta 6}
\begin{figure}[H]
\begin{lstlisting}
-- Consulta 6: pares de vendedores de la misma tienda, primero el más antiguo
SELECT a.nombre        AS vendedor_antiguo,
       a.fecha_ingreso AS ingreso_antiguo,
       b.nombre        AS vendedor_nuevo,
       b.fecha_ingreso AS ingreso_nuevo,
       t.nombre        AS tienda
FROM vendedores a
JOIN vendedores b
  ON b.id_tienda = a.id_tienda
 AND (a.fecha_ingreso < b.fecha_ingreso
      OR (a.fecha_ingreso = b.fecha_ingreso AND a.id_vendedor < b.id_vendedor))
JOIN tiendas t ON t.id_tienda = a.id_tienda
ORDER BY t.nombre, a.fecha_ingreso, b.fecha_ingreso;
\end{lstlisting}
\centering
\includegraphics[width=\textwidth]{imagenes/c06.png}
\caption{Resultado de la consulta 6}
\end{figure}


\clearpage
\section*{Consulta 7}
\begin{figure}[H]
\begin{lstlisting}
-- Consulta 7: monto en Montaña y Ruta por tienda, solo las que superan 2.000.000
SELECT t.nombre AS tienda,
       SUM(ve.unidades * b.precio) AS monto_montana_ruta
FROM tiendas t
JOIN vendedores v  ON v.id_tienda    = t.id_tienda
JOIN ventas ve     ON ve.id_vendedor = v.id_vendedor
JOIN bicicletas b  ON b.id_bicicleta = ve.id_bicicleta
WHERE b.categoria IN ('Montaña', 'Ruta')
GROUP BY t.id_tienda, t.nombre
HAVING SUM(ve.unidades * b.precio) > 2000000
ORDER BY monto_montana_ruta DESC;
\end{lstlisting}
\centering
\includegraphics[width=\textwidth]{imagenes/c07.png}
\caption{Resultado de la consulta 7}
\end{figure}


\clearpage
\section*{Consulta 8}
\begin{figure}[H]
\begin{lstlisting}
-- Consulta 8: bicicletas con precio sobre el promedio del catálogo
SELECT modelo, precio
FROM bicicletas
WHERE precio > (SELECT AVG(precio) FROM bicicletas)
ORDER BY precio DESC;
\end{lstlisting}
\centering
\includegraphics[width=0.58\textwidth]{imagenes/c08.png}
\caption{Resultado de la consulta 8}
\end{figure}

\begin{figure}[H]
\begin{lstlisting}
-- Consulta 8 (apoyo): valor del promedio del catálogo usado en la subconsulta
SELECT AVG(precio) AS promedio_catalogo
FROM bicicletas;
\end{lstlisting}
\centering
\includegraphics[width=0.58\textwidth]{imagenes/c09.png}
\caption{Promedio del catálogo usado en la subconsulta}
\end{figure}

\noindent\textbf{Explicación:} El promedio del catálogo es \$657.490; lo obtuve de la subconsulta \texttt{SELECT AVG(precio) FROM bicicletas}, la misma que usa el \texttt{WHERE}, ejecutada por separado en la segunda captura.

\clearpage
\section*{Consulta 9}
\begin{figure}[H]
\begin{lstlisting}
-- Consulta 9: vendedores con promedio de unidades por venta sobre el de la cadena
SELECT v.nombre,
       AVG(ve.unidades) AS promedio_unidades
FROM vendedores v
JOIN ventas ve ON ve.id_vendedor = v.id_vendedor
GROUP BY v.id_vendedor, v.nombre
HAVING AVG(ve.unidades) > (SELECT AVG(unidades) FROM ventas)
ORDER BY promedio_unidades DESC, v.nombre;
\end{lstlisting}
\centering
\includegraphics[width=\textwidth]{imagenes/c10.png}
\caption{Resultado de la consulta 9}
\end{figure}


\clearpage
\section*{Consulta 10}
\begin{figure}[H]
\begin{lstlisting}
-- Consulta 10: monto total vendido por cada vendedor (incluye los que no vendieron)
SELECT v.nombre,
       SUM(ve.unidades * b.precio) AS monto_total
FROM vendedores v
LEFT JOIN ventas ve    ON ve.id_vendedor = v.id_vendedor
LEFT JOIN bicicletas b ON b.id_bicicleta = ve.id_bicicleta
GROUP BY v.id_vendedor, v.nombre
ORDER BY monto_total DESC;
\end{lstlisting}
\centering
\includegraphics[width=\textwidth]{imagenes/c11.png}
\caption{Resultado de la consulta 10}
\end{figure}

\noindent\textbf{Explicación:} Fernanda Lillo y Rubén Castro aparecen con NULL porque no tienen ventas: el \texttt{LEFT JOIN} les deja \texttt{unidades} y \texttt{precio} en NULL y \texttt{SUM} sobre puros NULL devuelve NULL; para que mostraran cero habría que escribir \texttt{COALESCE(SUM(ve.unidades * b.precio), 0)}.

\end{document}
